\documentclass[11pt]{article}

% Shared notes settings for Elements of AIML lectures (UPES).
% Keep this file free of \documentclass to allow reuse via \input.

\usepackage[utf8]{inputenc}
\usepackage[T1]{fontenc}
\usepackage[a4paper,margin=1in]{geometry}
\usepackage{hyperref}
\usepackage{xurl}
\usepackage{booktabs}
\usepackage{amsmath}
\usepackage{amssymb}
\usepackage{listings}
\usepackage{xcolor}
\usepackage{enumitem}
\usepackage{parskip}

% Code listings (general-purpose).
\lstset{
  basicstyle=\ttfamily\small,
  keywordstyle=\color{blue},
  commentstyle=\color{gray},
  stringstyle=\color{teal},
  showstringspaces=false,
  columns=fullflexible,
  frame=single,
  framerule=0.2pt,
  breaklines=true
}

\setlist[itemize]{topsep=0.3em,itemsep=0.2em}
\setlist[enumerate]{topsep=0.3em,itemsep=0.2em}

% Simple per-section source line.
\newcommand{\notesources}[1]{\par\noindent{\small\color{gray}\textit{Sources:} \sloppy #1\par}}

% Unit-wise source strings for this subject (used by slides/notes).
% Path is relative to the lecture `latex/` folder (the compilation CWD).
\IfFileExists{../../../_shared/aiml-sources.tex}{%
  % Source strings for Elements of AIML (used by slides + notes).

\newcommand{\AIMLRepoURL}{https://github.com/tali7c/Elements-of-AIML}

\newcommand{\AIMLLabSources}{%
Provost and Fawcett, \textit{Data Science for Business}.;%
McKinney, \textit{Python for Data Analysis}.;%
WEKA Documentation (Waikato Environment for Knowledge Analysis), accessed 2026-02-28.;%
Matplotlib and Seaborn documentation, accessed 2026-02-28.%
}

\newcommand{\AIMLUnitISources}{%
Rich and Knight, \textit{Artificial Intelligence}, McGraw Hill, 2017.;%
Russell and Norvig, \textit{Artificial Intelligence: A Modern Approach} (4e), Ch. 1--2 (Introduction, Intelligent Agents).;%
Poole and Mackworth, \textit{Artificial Intelligence: Foundations of Computational Agents} (2e), selected sections.%
}

\newcommand{\AIMLUnitIISources}{%
Russell and Norvig, \textit{Artificial Intelligence: A Modern Approach} (4e), Ch. 7--9 (Logical Agents, First-Order Logic, Inference).;%
Huth and Ryan, \textit{Logic in Computer Science} (2e), selected sections on propositional and first-order logic.;%
Genesereth and Nilsson, \textit{Logical Foundations of Artificial Intelligence}, selected chapters.%
}

\newcommand{\AIMLUnitIIISources}{%
Mueller and Massaron, \textit{Machine Learning for Dummies}, For Dummies, 2016.;%
Mitchell, \textit{Machine Learning}, selected chapters on learning basics and evaluation.;%
Scikit-learn documentation, accessed 2026-02-28.%
}

\newcommand{\AIMLUnitIVSources}{%
Mueller and Massaron, \textit{Machine Learning for Dummies}, For Dummies, 2016.;%
James, Witten, Hastie, and Tibshirani, \textit{An Introduction to Statistical Learning} (2e), selected chapters.;%
Scikit-learn documentation, accessed 2026-02-28.%
}

\newcommand{\AIMLUnitVSources}{%
NIST, \textit{AI Risk Management Framework (AI RMF 1.0)}, 2023.;%
Barocas, Hardt, and Narayanan, \textit{Fairness and Machine Learning} (online book), selected sections.;%
Knaflic, \textit{Storytelling with Data}, selected chapters on communicating results.%
}
%
}{%
  % Faculty copies live under `private/` so their compile CWD is different.
  \IfFileExists{../../../../public/_shared/aiml-sources.tex}{%
    % Source strings for Elements of AIML (used by slides + notes).

\newcommand{\AIMLRepoURL}{https://github.com/tali7c/Elements-of-AIML}

\newcommand{\AIMLLabSources}{%
Provost and Fawcett, \textit{Data Science for Business}.;%
McKinney, \textit{Python for Data Analysis}.;%
WEKA Documentation (Waikato Environment for Knowledge Analysis), accessed 2026-02-28.;%
Matplotlib and Seaborn documentation, accessed 2026-02-28.%
}

\newcommand{\AIMLUnitISources}{%
Rich and Knight, \textit{Artificial Intelligence}, McGraw Hill, 2017.;%
Russell and Norvig, \textit{Artificial Intelligence: A Modern Approach} (4e), Ch. 1--2 (Introduction, Intelligent Agents).;%
Poole and Mackworth, \textit{Artificial Intelligence: Foundations of Computational Agents} (2e), selected sections.%
}

\newcommand{\AIMLUnitIISources}{%
Russell and Norvig, \textit{Artificial Intelligence: A Modern Approach} (4e), Ch. 7--9 (Logical Agents, First-Order Logic, Inference).;%
Huth and Ryan, \textit{Logic in Computer Science} (2e), selected sections on propositional and first-order logic.;%
Genesereth and Nilsson, \textit{Logical Foundations of Artificial Intelligence}, selected chapters.%
}

\newcommand{\AIMLUnitIIISources}{%
Mueller and Massaron, \textit{Machine Learning for Dummies}, For Dummies, 2016.;%
Mitchell, \textit{Machine Learning}, selected chapters on learning basics and evaluation.;%
Scikit-learn documentation, accessed 2026-02-28.%
}

\newcommand{\AIMLUnitIVSources}{%
Mueller and Massaron, \textit{Machine Learning for Dummies}, For Dummies, 2016.;%
James, Witten, Hastie, and Tibshirani, \textit{An Introduction to Statistical Learning} (2e), selected chapters.;%
Scikit-learn documentation, accessed 2026-02-28.%
}

\newcommand{\AIMLUnitVSources}{%
NIST, \textit{AI Risk Management Framework (AI RMF 1.0)}, 2023.;%
Barocas, Hardt, and Narayanan, \textit{Fairness and Machine Learning} (online book), selected sections.;%
Knaflic, \textit{Storytelling with Data}, selected chapters on communicating results.%
}
%
  }{%
    % Source strings for Elements of AIML (used by slides + notes).

\newcommand{\AIMLRepoURL}{https://github.com/tali7c/Elements-of-AIML}

\newcommand{\AIMLLabSources}{%
Provost and Fawcett, \textit{Data Science for Business}.;%
McKinney, \textit{Python for Data Analysis}.;%
WEKA Documentation (Waikato Environment for Knowledge Analysis), accessed 2026-02-28.;%
Matplotlib and Seaborn documentation, accessed 2026-02-28.%
}

\newcommand{\AIMLUnitISources}{%
Rich and Knight, \textit{Artificial Intelligence}, McGraw Hill, 2017.;%
Russell and Norvig, \textit{Artificial Intelligence: A Modern Approach} (4e), Ch. 1--2 (Introduction, Intelligent Agents).;%
Poole and Mackworth, \textit{Artificial Intelligence: Foundations of Computational Agents} (2e), selected sections.%
}

\newcommand{\AIMLUnitIISources}{%
Russell and Norvig, \textit{Artificial Intelligence: A Modern Approach} (4e), Ch. 7--9 (Logical Agents, First-Order Logic, Inference).;%
Huth and Ryan, \textit{Logic in Computer Science} (2e), selected sections on propositional and first-order logic.;%
Genesereth and Nilsson, \textit{Logical Foundations of Artificial Intelligence}, selected chapters.%
}

\newcommand{\AIMLUnitIIISources}{%
Mueller and Massaron, \textit{Machine Learning for Dummies}, For Dummies, 2016.;%
Mitchell, \textit{Machine Learning}, selected chapters on learning basics and evaluation.;%
Scikit-learn documentation, accessed 2026-02-28.%
}

\newcommand{\AIMLUnitIVSources}{%
Mueller and Massaron, \textit{Machine Learning for Dummies}, For Dummies, 2016.;%
James, Witten, Hastie, and Tibshirani, \textit{An Introduction to Statistical Learning} (2e), selected chapters.;%
Scikit-learn documentation, accessed 2026-02-28.%
}

\newcommand{\AIMLUnitVSources}{%
NIST, \textit{AI Risk Management Framework (AI RMF 1.0)}, 2023.;%
Barocas, Hardt, and Narayanan, \textit{Fairness and Machine Learning} (online book), selected sections.;%
Knaflic, \textit{Storytelling with Data}, selected chapters on communicating results.%
}
%
  }%
}


\title{Elements of AIML\\Unit 05 -- Lecture 02 Notes\\ML Applications in Banking, Security \& Healthcare}
\author{Tofik Ali}
\date{\today}

\begin{document}
\maketitle
\tableofcontents

\section{Lecture Overview}
This lecture uses three domains (banking, security, healthcare) to show:
\begin{itemize}[leftmargin=*]
  \item how ML tasks map to real products,
  \item why evaluation is domain-specific,
  \item how drift, privacy, and safety constraints change deployment decisions.
\end{itemize}
\notesources{\AIMLUnitVSources}

\section{How to use these notes (self-study)}
Treat each domain as a case study.
For each one, identify: (1) the ML task type (classification/regression/anomaly), (2) the most important metric (based on cost of errors), and (3) the main deployment risks (drift, adversaries, privacy, safety).
This is the same checklist you will use in projects and interviews.

\section{Learning Outcomes}
\begin{itemize}[leftmargin=*]
  \item List common ML tasks in banking, security, and healthcare.
  \item Identify common data sources and labeling challenges.
  \item Explain key risks: imbalance, adversaries, drift, privacy, and safety.
  \item Describe monitoring needs for deployed models.
\end{itemize}

\section{Keywords}
\begin{itemize}[leftmargin=*]
  \item \textbf{Fraud detection:} rare-event classification.
  \item \textbf{Anomaly detection:} detect unusual behavior relative to normal patterns.
  \item \textbf{Concept drift:} changing data patterns over time.
  \item \textbf{False alarm fatigue:} too many false positives reduce trust and response quality.
  \item \textbf{Clinical validation:} testing model effectiveness in real clinical workflow.
\end{itemize}

\section{Abbreviations and key terms}
\begin{itemize}[leftmargin=*]
  \item \textbf{AML}: Anti-Money Laundering; \textbf{KYC}: Know Your Customer.
  \item \textbf{PII}: personally identifiable information.
  \item \textbf{EHR}: electronic health record; \textbf{PHI}: protected health information.
  \item \textbf{Drift}: change in data/label patterns over time; \textbf{monitoring}: ongoing checks in production.
\end{itemize}

\section{Banking applications}
\subsection{Fraud detection}
Fraud is usually rare, so datasets are imbalanced. Key considerations:
\begin{itemize}[leftmargin=*]
  \item Use precision/recall and cost-based evaluation.
  \item False positives annoy customers; false negatives lose money.
  \item Drift occurs as fraudsters adapt.
\end{itemize}

\subsection{Typical data sources}
\begin{itemize}[leftmargin=*]
  \item transaction history (amount, merchant category, time, location),
  \item device/browser signals (fraud often correlates with device patterns),
  \item customer/account history (age of account, past chargebacks).
\end{itemize}
Labels can be delayed (fraud confirmed after investigation), so training/evaluation pipelines must handle late-arriving ground truth.

\subsection{Credit scoring}
Credit scoring predicts default risk. Consider:
\begin{itemize}[leftmargin=*]
  \item fairness and regulatory constraints,
  \item explainability (users may request reasons),
  \item monitoring for economic shifts.
\end{itemize}

\subsection{What can go wrong (banking)}
\begin{itemize}[leftmargin=*]
  \item \textbf{Proxy bias}: using features that strongly correlate with protected attributes.
  \item \textbf{Policy mismatch}: model optimized for accuracy but business cost is different (e.g., false declines are very costly).
  \item \textbf{Drift}: sudden economic changes can invalidate older patterns.
\end{itemize}

\section{Security applications}
\subsection{Spam/phishing detection}
Classic classification problem with adversarial adaptation.
Attackers may change wording and links to bypass detection.

\subsection{Intrusion/anomaly detection}
Often framed as anomaly detection:
\begin{itemize}[leftmargin=*]
  \item learn patterns of normal behavior,
  \item flag deviations as suspicious.
\end{itemize}
Trade-off: too sensitive $\rightarrow$ too many alerts; too insensitive $\rightarrow$ missed attacks.

\subsection{What makes security different?}
\begin{itemize}[leftmargin=*]
  \item \textbf{Adversaries}: attackers intentionally try to evade detection (adaptive behavior).
  \item \textbf{Feedback loops}: once defenses change, attacker behavior shifts.
  \item \textbf{Operational constraints}: security teams must handle alerts; too many FP can cause ``alert fatigue''.
\end{itemize}

\section{Healthcare applications}
\subsection{Risk prediction and decision support}
Predicting risk (e.g., readmission) can support clinicians. Key issues:
\begin{itemize}[leftmargin=*]
  \item privacy and consent,
  \item careful validation to avoid harm,
  \item transparency and workflow integration.
\end{itemize}

\subsection{Data challenges (healthcare)}
\begin{itemize}[leftmargin=*]
  \item missing/noisy measurements (not every test is done for every patient),
  \item label ambiguity (diagnoses can be revised),
  \item distribution shift across hospitals, devices, and populations.
\end{itemize}

\subsection{Imaging assistance}
Models may classify or segment images. Evaluation must consider:
\begin{itemize}[leftmargin=*]
  \item sensitivity (recall) for critical conditions,
  \item false negatives can be dangerous,
  \item dataset shift across hospitals/devices.
\end{itemize}

\subsection{Clinical validation (why ``offline accuracy'' is not enough)}
Healthcare models should be validated in realistic settings:
\begin{itemize}[leftmargin=*]
  \item prospective/retrospective validation with appropriate splits,
  \item evaluation of workflow impact (time saved, reduced errors),
  \item human factors (interpretability, trust, escalation).
\end{itemize}

\section{Concept drift and monitoring}
Across these domains, drift is common:
\begin{itemize}[leftmargin=*]
  \item fraud patterns change,
  \item malware evolves,
  \item clinical populations and protocols change.
\end{itemize}
Monitoring includes:
\begin{itemize}[leftmargin=*]
  \item tracking data statistics (feature distributions),
  \item tracking performance metrics (when labels are available),
  \item alerting and rollback/retraining strategies.
\end{itemize}

\section{Practice (with answers)}
\subsection*{P1. Give one monitoring signal that does not require labels.}
\textbf{Answer:} feature distribution shift (e.g., missing rate, mean/std), or model score distribution changes.

\subsection*{P2. Why does ``false alarm fatigue'' happen in security systems?}
\textbf{Answer:} if a system generates too many false positives, humans stop trusting alerts and may ignore real incidents.

\section{Common pitfalls (and how to avoid them)}
\begin{itemize}[leftmargin=*]
  \item \textbf{Using the wrong metric:} fraud and intrusion are imbalanced; accuracy is rarely enough.
  \item \textbf{Ignoring delayed labels:} in fraud/healthcare, ground truth can arrive late; plan evaluation accordingly.
  \item \textbf{Assuming threats are stationary:} in security, attackers adapt; drift is adversarial and continuous.
  \item \textbf{Deploying without workflow fit:} in healthcare, a model must be validated and usable in the clinical workflow.
  \item \textbf{Neglecting privacy constraints:} banking/healthcare data often includes PII/PHI; access and governance are mandatory.
\end{itemize}

\section{Exit questions (with answers)}

\subsection*{Q1. Why is class imbalance common in fraud detection?}
\textbf{Answer:} fraud events are rare compared to normal transactions, so positive class examples are much fewer.

\subsection*{Q2. One example of drift in security.}
\textbf{Answer:} attackers change phishing email wording, domains, or tactics; malware signatures evolve; network behavior changes after new software updates.

\subsection*{Q3. In healthcare, why can false negatives be more costly?}
\textbf{Answer:} missing a critical condition can directly lead to harm, delayed treatment, or death; therefore recall is often prioritized.

\subsection*{Q4. Two monitoring signals after deployment.}
\textbf{Answer:} distribution shift in key features; false positive rate; model score distribution; latency; and periodic ground-truth performance once labels arrive.

\section{Summary}
Applications differ by constraints, but successful deployment requires correct metrics, robust evaluation, drift monitoring, and privacy/safety controls.
\notesources{\AIMLUnitVSources}

\end{document}
